\section{结论}
\projectname{} 适合作为 DP Technology / Deep Potential 实习申请中的 demo proposal，原因在于它同时满足四个条件：第一，问题具有明确的科学意义，直接对应实验谱图到分子结构反演这一高价值逆问题；第二，方案具备最小可执行路径，可以先交付一个编码--候选--前向一致性重排的闭环原型；第三，方法论与 AI for Science 高度契合，强调图结构、几何约束、物理一致性与可追踪验证；第四，体系具有自然扩展性，后续可以平滑过渡到 graph diffusion、pairwise distance 与 TFN-Transformer 等更强的几何深度学习模块。

从实习 demo 的角度看，这种“先做可运行闭环，再逐步补强科学深度”的策略比直接承诺一个完整大模型更可信，也更能体现对问题本质和工程边界的理解。若后续补齐真实图示、最小数据流和样例输出，本提案即可作为一份兼具研究视角、技术路线和执行可行性的申请材料。
